% Generado: no editar cifras manualmente.
\begin{table}[!htbp]
\centering
\caption{Sexo registrado por ronda}\label{tab:cp1-sexo}
\small\setlength{\tabcolsep}{4pt}
\renewcommand{\arraystretch}{1.15}
\begin{tabular}{p{0.28\linewidth}rrrrrr}
\toprule
\textbf{Categoría} & \textbf{2001} & \textbf{2003} & \textbf{2012} & \textbf{2015} & \textbf{2018} & \textbf{2021} \\
\midrule
Hombre & \shortstack[r]{6\,503\\(42,82\%)} & \shortstack[r]{6\,108\\(42,86\%)} & \shortstack[r]{8\,000\\(43,33\%)} & \shortstack[r]{6\,739\\(42,15\%)} & \shortstack[r]{7\,761\\(42,53\%)} & \shortstack[r]{7\,446\\(42,46\%)} \\
Mujer & \shortstack[r]{8\,683\\(57,18\%)} & \shortstack[r]{8\,142\\(57,14\%)} & \shortstack[r]{10\,465\\(56,67\%)} & \shortstack[r]{9\,249\\(57,85\%)} & \shortstack[r]{10\,488\\(57,47\%)} & \shortstack[r]{10\,092\\(57,54\%)} \\
\bottomrule
\end{tabular}
\par\medskip
\parbox{\linewidth}{\footnotesize Fuente: elaboración propia con agregados de CP1, ENASEMsimple v2 (2001--2021). Frecuencias sin ponderación. Cada celda: n y (\%). Denominador: todos los registros de la ronda de la tabla de población, incluidos faltantes. Las categorías suman N; un código ausente tiene frecuencia cero.}
\end{table}
